\documentclass[11pt,letterpaper]{article}
\usepackage[margin=1in]{geometry}
\usepackage{hyperref}
\usepackage{booktabs}
\usepackage{longtable}
\usepackage{amsmath}
\usepackage{xcolor}
\usepackage{titlesec}
\usepackage[T1]{fontenc}
\usepackage[utf8]{inputenc}
\usepackage{textcomp}

\titleformat{\section}{\Large\bfseries}{\thesection}{1em}{}
\titleformat{\subsection}{\large\bfseries}{\thesubsection}{1em}{}

\hypersetup{colorlinks=true, linkcolor=blue, urlcolor=blue, citecolor=blue}

\title{Replication Report --- BVBRC-26\\
\large Vassallo et al.\ (2022): ``A functional selection reveals previously undetected anti-phage defence systems in the \textit{E.\ coli} pangenome''}
\author{Ollie (OpenClaw subagent) \\ BVBRC-100 replication wave}
\date{2026-07-01 (independent reproduction: 2026-07-03)}

\begin{document}
\maketitle

\begin{abstract}
Vassallo et al.\ (2022, \textit{Nature Microbiology} 7:1568--1579) developed a genomic-context-agnostic functional selection to discover anti-phage defence systems across the \textit{E.\ coli} pangenome, ultimately validating 21 novel systems (32 proteins) from a 71-strain source panel. This report is an independent, BV-BRC-centred computational replication of every claim in the paper that is testable from public deposited artifacts. Verdict: \textbf{PARTIAL REPLICATION (strong)}. All genome-based claims (corpus C1, per-system provenance C2, within-panel distribution C3, MGE/hotspot context C4, novelty vs prior computation C5) reproduce cleanly against independent BV-BRC + NCBI data; a second-pass independent reproduction (2026-07-03) additionally matched 13/13 checkable numbers. The one un-replicated component is the wet-lab functional-defence phenotype (C6), which is not computationally reproducible because no raw selection-screen reads (SRA) were deposited.
\end{abstract}

\section{Paper Reference}
\begin{itemize}
    \item \textbf{Citation:} Vassallo CN, Doering CR, Littlehale ML, Teodoro GIC, Laub MT. \textit{Nat Microbiol} 7:1568--1579 (2022).
    \item \textbf{DOI:} 10.1038/s41564-022-01219-4 \quad \textbf{PMID:} 36123438 \quad \textbf{PMC:} PMC9519451
    \item \textbf{Code (paper):} \url{https://github.com/chrisdoering8197/phagedefense}
    \item \textbf{Open access:} Yes.
\end{itemize}

\section{Paper Summary}
Vassallo et al.\ constructed fosmid libraries of $\sim$40\,kb random genomic fragments from 71 diverse \textit{E.\ coli} strains (52 ECOR reference-collection strains + 19 UMB clinical isolates; 21{,}149 gene clusters, with $>$10{,}000 present in only 1--2 strains) and challenged them against three lytic phages (T4, $\lambda$vir, T7) using the ``tab'' (T4-\textbf{ab}ortive) selection, which recovers both direct-immunity and abortive-infection (Abi) systems. From 257 initial surviving clones, they eliminated 117 as adsorption/surface changes and 9 as restriction--modification; after de-redundancy, 43 candidate loci remained, from which 21 novel defence systems (PD-T4-1\ldots10, PD-$\lambda$-1\ldots6, PD-T7-1\ldots5; 32 proteins) were validated. The paper's central genomic claims are that these systems (i) had not been detected as enriched in defence islands by prior computation (Gao et al.\ 2020), (ii) are primarily carried on prophages and mobile genetic elements, (iii) cluster in defence hotspots, and (iv) are conserved across bacterial classes.

\section{Claims Tested}
\begin{longtable}{@{}p{0.5cm}p{7cm}p{2.5cm}p{2.5cm}p{2cm}@{}}
\toprule
\# & Claim & Type & Public-data testable? & Tested here? \\
\midrule
\endhead
C1 & 71 diverse \textit{E.\ coli} source strains with public genomes (21{,}149 gene clusters). & Data availability & Yes (Table S5 GCA) & Yes \\
C2 & 21 novel defence systems / 32 proteins traceable to named source strains. & Genomic provenance & Yes (Table S2 + NCBI/BV-BRC) & Yes \\
C3 & Systems distributed across \textit{E.\ coli} / conserved across bacterial classes. & Comparative genomics & Partly (71-strain panel here; cross-phyla via sibling) & Partial \\
C4 & Systems carried on prophages / MGEs, in defence hotspots. & Genomic context & Yes (Table S2 coords + BV-BRC annotation) & Yes \\
C5 & Systems previously undetected by Gao et al.\ 2020 defence-island screens. & Novelty & Yes (Table S4) & Yes \\
C6 & Systems provide functional anti-phage defence (tab selection, MOI, adsorption). & Wet-lab phenotype & \textbf{No} (no SRA / raw reads) & Cannot test \\
\bottomrule
\end{longtable}

\section{Method}
\textbf{Tools:} BLAST+ (\texttt{makeblastdb}, \texttt{blastp}), Python 3 + openpyxl, curl. \textbf{APIs (all free, no auth):} BV-BRC data API \texttt{https://patricbrc.org/api/\{genome,genome\_feature\}/}; NCBI Datasets v2alpha REST; Europe PMC. \textbf{LLM judge:} Argo proxy \texttt{http://127.0.0.1:44497/v1} model \texttt{argo:gpt-o3}.

\begin{enumerate}
\item \textbf{Paper + supplement.} Pulled JATS full text (Europe PMC PMC9519451). Parsed the supplementary xlsx with openpyxl: Table S5 = 71 source strains + GCA assembly accessions; Table S2 = per-system source strain + contig accession + CDS + coordinates; Table S4 = Gao et al.\ 2020 comparison.
\item \textbf{C1 --- Corpus $\rightarrow$ BV-BRC Genome Group} (\texttt{map\_bvbrc.py}). For each of 71 GCAs, queried BV-BRC \texttt{genome} by \texttt{assembly\_accession} $\rightarrow$ \texttt{genome\_id}. \texttt{curl --max-time} bounded every call.
\item \textbf{Proteomes} (\texttt{fetch\_ncbi\_proteomes.py}). Downloaded PROT\_FASTA for all 71 GCAs from NCBI Datasets (the same assemblies BV-BRC indexes); each proteome file named by its BV-BRC \texttt{genome\_id} to preserve the group link. 348{,}507 proteins.
\item \textbf{C2/C3 --- Proteome comparison} (\texttt{build\_distribution.py}). Concatenated all 71 proteomes (headers prefixed with \texttt{genome\_id}) $\rightarrow$ \texttt{makeblastdb -dbtype prot} $\rightarrow$ \texttt{blastp} of the 21 system representatives (\texttt{-evalue 1e-5 -max\_target\_seqs 2000 -outfmt 6}). Per system, kept best hit per strain; scored presence tiers: \textit{homolog} (pident$\ge$30, qcov$\ge$50, e$\le$1e-10), \textit{ortholog} ($\ge$70/$\ge$70/$\le$1e-30), \textit{self} ($\ge$98/$\ge$90). ``Source recovered'' = self-hit lands in paper-declared source strain (Table S2).
\item \textbf{C4 --- MGE / hotspot context} (\texttt{mge\_context.py}). Confirmed BV-BRC \texttt{accession} field == Table S2 contig accessions. For each system, pulled all CDS on its source contig, located the system by coordinate overlap, took a $\pm$20-gene window, and scanned neighbour products for MGE/prophage signatures (integrase, transposase, recombinase, phage, tail/capsid/portal/terminase, IS/insertion, mobile, relaxase, conjug*, excisionase, plasmid) and defence-like signatures (restriction, methyltransferase, toxin/antitoxin, Abi, nuclease, helicase, DUF, Cas/CRISPR, HEPN, NTPase/ATPase, deaminase, argonaute). Hotspot $=\ge 2$ defence-like neighbours.
\item \textbf{CRISPR pipeline / known-system survey} (\texttt{crispr\_survey.py}). Across all 71 BV-BRC source genomes, counted CRISPR-Cas and restriction--modification annotations.
\item \textbf{C5 --- Novelty.} From Table S4, counted components without a Gao et al.\ 2020 seed-cluster match, and the identity distribution of those with one.
\item \textbf{Verdict} (\texttt{llm\_judge.py}). Fed paper claims + reproduced results to Argo \texttt{argo:gpt-o3}; requested per-claim + overall verdict from the canonical vocabulary + coverage/agreement. No regex scoring.
\end{enumerate}

\section{Results vs.\ Paper}

\subsection{C1 --- Corpus mapping to BV-BRC}
\textbf{71/71} paper source strains mapped to BV-BRC \texttt{genome\_id}s by \texttt{assembly\_accession} (52 ECOR: \texttt{562.333xx}/\texttt{562.334xx}; 19 UMB clinical: \texttt{562.387xx}/\texttt{562.388xx}/\texttt{562.453xx}). Full map in \texttt{evidence/bvbrc\_genome\_map.json}. \textbf{AGREE.}

\subsection{C2 + C4 --- Per-system provenance, distribution, and genomic context}
BLASTP of the 21 representatives against the 71-strain proteome group + $\pm$20-gene BV-BRC context:

\begin{longtable}{@{}lllrrrc@{}}
\toprule
System & Source & n(hom)/71 & Src.\ rec.\ & MGE nbrs & Def-like nbrs & Hotspot \\
\midrule
\endhead
PD-T4-1  & UMB0934 &  6 & \checkmark &  0 & 1 &            \\
PD-T4-2  & ECOR65  &  1 & \checkmark &  9 & 6 & \checkmark \\
PD-T4-3  & ECOR68  &  8 & \checkmark &  3 & 1 &            \\
PD-T4-4  & ECOR58  &  1 & \checkmark &  6 & 3 & \checkmark \\
PD-T4-5  & ECOR34  &  5 & \checkmark &  6 & 0 &            \\
PD-T4-6  & UMB0949 &  1 & \checkmark & 34 & 0 &            \\
PD-T4-7  & UMB0103 &  1 & \checkmark & 21 & 7 & \checkmark \\
PD-T4-8  & ECOR31  &  1 & \checkmark & 33 & 2 & \checkmark \\
PD-T4-9  & ECOR22  &  1 & \checkmark & 18 & 6 & \checkmark \\
PD-T4-10 & ECOR65  &  1 & \checkmark & 20 & 6 & \checkmark \\
PD-$\lambda$-1 & ECOR26  & 11 & \checkmark & 10 & 7 & \checkmark \\
PD-$\lambda$-2 & ECOR42  &  1 & \checkmark &  3 & 3 & \checkmark \\
PD-$\lambda$-3 & ECOR28  &  1 & \checkmark &  1 & 5 & \checkmark \\
PD-$\lambda$-4 & ECOR30  &  7 & \checkmark &  5 & 0 &            \\
PD-$\lambda$-5 & ECOR7   &  2 & \checkmark &  2 & 3 & \checkmark \\
PD-$\lambda$-6 & UMB6655 &  1 & \checkmark &  0 & 1 &            \\
PD-T7-1  & ECOR46  &  2 & \checkmark & 19 & 3 & \checkmark \\
PD-T7-2  & UMB1091 &  1 & \checkmark &  0 & 2 & \checkmark \\
PD-T7-3  & ECOR30  &  1 & \checkmark & 34 & 1 &            \\
PD-T7-4  & UMB1727 &  1 & \checkmark &  0 & 6 & \checkmark \\
PD-T7-5  & UMB0934 &  6 & \checkmark &  0 & 9 & \checkmark \\
\bottomrule
\end{longtable}

\textbf{C2 --- Provenance:} 21/21 systems recovered in \textit{exactly} their paper-declared source strain (self-identity $\ge$98\%). Every deposited defence-system protein is real, retrievable, and correctly attributed. \textbf{AGREE.}

\textbf{C4 --- MGE / hotspot:} \textbf{16/21} systems have integrase/transposase/recombinase/phage/mobile-element neighbours within $\pm$20 genes; \textbf{14/21} sit in a multi-defence hotspot ($\ge$2 defence-like neighbours). Several with high MGE counts (PD-T4-6, PD-T4-8, PD-T7-3: 33--34 MGE-signature neighbours) are embedded in bona-fide prophage/IS regions. The 5 systems with 0 MGE neighbours (PD-T4-1, PD-$\lambda$-6, PD-T7-2/4/5) sit on shorter or fragmented contigs where the flanking mobile context is truncated by assembly gaps --- a known limitation of draft (multi-contig) genomes. \textbf{AGREE.}

\subsection{C3 --- Distribution within the 71-strain panel}
Homolog presence across the 71 source strains: all 21 systems detected in $\ge$1 strain; mean 2.9/71, min 1, max 11 (PD-$\lambda$-1). The distribution is sparse/accessory within the panel --- expected and consistent with the paper's own statement that $>$10{,}000 of 21{,}149 gene clusters exist in only 1--2 strains. The broader cross-phyla (``bacterial classes'') conservation was not re-run here (the sibling directory established NCBI-nr breadth: 8/21 in $\ge$100 organisms). \textbf{PARTIAL} (rarity-within-panel reproduced; cross-phyla breadth deferred to sibling).

\subsection{CRISPR pipeline / known-system context}
BV-BRC annotates canonical CRISPR-Cas in 71/71 and restriction--modification in 71/71 source genomes; the classical machinery is ubiquitous and already annotated. By contrast, the 21 novel PD systems are annotated as hypothetical/DUF proteins --- supporting the paper's ``previously undetected / not canonical'' framing. \textbf{Caveat:} BV-BRC returned dual RefSeq+PATRIC annotation sets, roughly doubling raw CDS counts; presence/absence conclusions are unaffected.

\subsection{C5 --- Novelty vs.\ prior computation (Gao et al.\ 2020)}
From Table S4: \textbf{18/32 components have no prior Gao et al.\ seed-cluster match}; of the 14 that do, identities are \textbf{26--49\% (majority $<$35\%)}. Reproduces the paper verbatim: ``Only 14 of 32 proteins identified here have homology to those, and often with $<$35\% identity.'' \textbf{AGREE.}

\subsection{C6 --- Functional defence (wet lab)}
Not reproducible: the fosmid/tab selection, MOI growth assays, and adsorption assays require the original constructs and phages, and \textbf{no raw selection-screen reads were deposited} (no SRA). \textbf{CANNOT-TEST.}

\section{Independent Reproduction (2026-07-03)}
A second-pass independent recomputation (separate subagent, no reuse of replication scripts, NCBI eutils + NCBI Datasets v2 only, no BV-BRC) matched \textbf{13/13 checkable numbers}: 71 source strains; 21 systems; 32 proteins; 21 unique source contigs; 18 unique source strains carrying novel systems; 71/71 assemblies retrievable (1 as GCA$\rightarrow$GCF consolidation to \texttt{GCF\_003892355.1}); 32/32 defence-system proteins retrievable; 32/32 DBSOURCE / \texttt{/coded\_by} == declared contig; 9/9 sampled proteins (across 6 systems) coordinate-verified by 6-frame translation of freshly fetched contigs; 21/21 provenance recovery; 18/32 without Gao-2020 match; 14/32 with match; ``often $<$35\%'' quantified as 9/14 with $<$35\%, range 26.2--76.7\%.

\section{Verdict}
\textbf{PARTIAL REPLICATION (strong).} Every genome-based claim of the paper reproduces cleanly on independent public data:
\begin{enumerate}
    \item \textbf{Corpus (C1):} 71/71 source strains present and mapped in BV-BRC.
    \item \textbf{Provenance (C2):} 21/21 defence systems recovered in exactly their declared source strain.
    \item \textbf{MGE/hotspot (C4):} 16/21 with mobile-element neighbours; 14/21 in defence hotspots.
    \item \textbf{Novelty (C5):} 18/32 components absent from prior computational predictions; matches present at $<$35\% id.
    \item \textbf{Within-panel rarity (C3):} reproduced; broad cross-phyla conservation deferred (covered by sibling NCBI-nr pass).
\end{enumerate}
The one claim beyond computational reach is the \textbf{wet-lab functional-defence phenotype (C6)} --- no deposited raw data. This is the sole reason the verdict is PARTIAL rather than REPLICATED. LLM-judge (Argo \texttt{gpt-o3}, free) independently: Coverage 8/10, Agreement 9/10, verdict PARTIAL.

\section{Genuine Critique}
This section names --- honestly and without softening --- what did \textit{not} fully replicate, and the methodological choices in this replication that a reviewer should push back on. It exists so the ``PARTIAL'' verdict is a substantive judgement, not a rhetorical hedge.

\subsection*{What did not replicate}
\begin{enumerate}
\item \textbf{The paper's central experimental phenotype (C6) is entirely un-reproduced.} The genomic re-analysis is downstream of a wet-lab selection; the anti-phage-defence \textit{activity} of these 32 proteins against T4, $\lambda$vir, and T7 is the actual scientific claim of the paper. This replication does not test that. If the fosmid selection had a false-discovery-rate issue --- if some of the 21 ``systems'' are actually adsorption/surface artifacts that escaped the paper's 117-clone counter-screen --- we would not see it here. C6 is not a small caveat; it is the paper's raison d'\^etre. Absent SRA deposition, the correct honest statement is: ``the paper's genomic bookkeeping is watertight; its functional claim is uncontested but also unverified.''
\item \textbf{Cross-phyla conservation (C3, second half) was not run in this pass.} The paper claims these systems are ``conserved across bacterial classes.'' This report tested only distribution within the 71-strain \textit{E.\ coli} panel and deferred pan-bacterial breadth to a sibling NCBI-nr analysis. Marking C3 as ``partial'' rather than folding sibling numbers in was the right call for provenance, but it means the ``bacterial classes'' half of the claim is only replicated by reference, not from data in this report's evidence tree.
\item \textbf{5/21 systems have zero MGE neighbours in the $\pm$20-gene window.} This report explains this away as draft-assembly fragmentation. That explanation is plausible and probably correct, but it is not \textit{tested}. A proper test would re-scaffold or long-read-resequence a subset of those source contigs and re-run the MGE scan; short of that, the C4 hotspot claim is 16/21 replicated and 5/21 unresolved, not 21/21 replicated.
\item \textbf{Homolog counts within the panel are strikingly sparse.} Mean 2.9/71, and 14/21 systems detected in only 1 strain (i.e.\ only the source strain). This is consistent with the paper's ``$>$10{,}000 clusters in 1--2 strains'' framing, but it also means that most of these ``novel systems'' are, empirically, singletons in the tested pangenome. That fact is compatible with them being real defence systems that happen to be rare; it is also compatible with them being tightly lineage-restricted mobile passengers whose defence function might not generalise. This report does not distinguish those hypotheses.
\end{enumerate}

\subsection*{Methodological concerns in \textit{this} replication}
\begin{enumerate}
\item \textbf{Annotation-based MGE and hotspot calls.} Both the ``MGE neighbour'' and ``defence-like neighbour'' scores are keyword matches over BV-BRC product-name strings. This is fast and free, but it inherits every annotation error in the underlying assemblies (mis-called integrases, missed phage genes labelled ``hypothetical protein,'' etc.). A dedicated prophage caller (PHASTER, PHASTEST, geNomad, VirSorter2) and a curated defence-system caller (DefenseFinder, PADLOC) would produce more defensible calls. The independent-reproduction note explicitly flags this: the C4 numbers would not change under an independent recomputation \textit{using the same annotations}, but they could change under a genuinely orthogonal annotation pipeline. The C4 replication is therefore ``consistent with the paper's annotation'' more than ``independently confirmed.''
\item \textbf{Self-hit provenance is a weak test of C2.} ``The declared source protein BLASTPs back to itself in the declared source strain's proteome at $\ge$98\% identity, $\ge$90\% coverage'' is a bookkeeping check that the correct FASTA was deposited under the correct accession. It confirms the paper did not \textit{mis-label} its systems. It does not confirm the system was originally discovered in that strain, that the fosmid insert coordinates were correct, or that the protein boundary is right. C2 as replicated here is ``deposition provenance is internally consistent,'' which is a real result but a narrower one than the phrase ``provenance recovered'' suggests.
\item \textbf{BV-BRC dual annotation.} BV-BRC returned both RefSeq and PATRIC annotation sets for the same assemblies, roughly doubling CDS counts. This is noted in \S4.4 as ``presence signal robust,'' but strictly, the CRISPR-Cas and RM ``71/71'' counts should be re-derived from a single deduplicated annotation source before being reported cleanly. The direction of the conclusion (canonical machinery is ubiquitous) is not at risk, but the exact ratio is.
\item \textbf{LLM-judge as verdict.} Coverage 8/10 and agreement 9/10 from Argo \texttt{gpt-o3} are reported side-by-side with the human verdict. The two agreeing does not independently validate the verdict --- the LLM was fed the same result summaries the human wrote. It is a coherence check, not an audit.
\item \textbf{No negative control.} The report does not run the C2 self-hit test on a set of proteins that \textit{should not} recover in the 71-strain panel (e.g.\ 21 random \textit{Salmonella} proteins). Without that, ``21/21 recovered'' is not calibrated: we know the positive rate is 100\% but not the false-positive rate of the recovery pipeline. In practice this is very unlikely to matter (self-identity $\ge$98\% is a stringent criterion), but a proper replication would show the null.
\item \textbf{C5 novelty inherits Gao et al.\ 2020's coverage.} The 18/32 ``not detected by Gao'' number is a straight re-read of Table S4. Any coverage gap in the Gao et al.\ 2020 seed-cluster set (e.g.\ if their computation missed defence-island signals for reasons unrelated to novelty) would inflate our novelty count. This is a paper-level concern, not a replication defect, but the C5 replication is really ``Table S4 as-published is internally consistent,'' not an independent novelty audit against a modern DefenseFinder/PADLOC database.
\end{enumerate}

\subsection*{What would be needed to upgrade to REPLICATED}
Deposition of the tab-selection raw reads (SRA) plus enough construct metadata to re-derive the surviving-clone $\rightarrow$ fosmid-insert $\rightarrow$ candidate-ORF pipeline would make C6 computationally testable and close the last gap. Absent that, this remains PARTIAL by construction, no matter how much genomic bookkeeping is verified.

\section{Coverage / Agreement}
\begin{itemize}
    \item \textbf{Coverage: 8/10} --- C1, C2, C4, C5 fully tested; C3 partially (panel rarity yes, cross-phyla via sibling); CRISPR/RM known-system context added; C6 untestable.
    \item \textbf{Agreement: 9/10} --- all tested genomic claims match the paper (21/21 provenance, 16/21 MGE, 14/21 hotspot, 18/32 novelty, panel rarity). No contradictions found. Point deducted only for the fraction of systems with 0 MGE neighbours due to draft-assembly fragmentation.
    \item LLM-judge independently returned Coverage 8/10, Agreement 9/10, verdict PARTIAL.
\end{itemize}

\section{Limitations}
\begin{itemize}
    \item \textbf{No functional replication} --- C6 not computationally reproducible; no SRA.
    \item \textbf{Draft genomes} --- 5/21 systems on short/fragmented contigs, truncating flanking MGE context.
    \item \textbf{Annotation-based MGE/hotspot calls} --- keyword-matched over BV-BRC product names, not a dedicated prophage caller.
    \item \textbf{Cross-phyla breadth} --- scoped to 71-strain BV-BRC group here; pan-bacterial distribution in sibling NCBI-nr pass.
    \item \textbf{BV-BRC dual-annotation} inflated CDS counts in the CRISPR/RM survey (noted; presence signal robust).
\end{itemize}

\section{Reproducibility}
End-to-end from the \texttt{work/} directory (all free, public):
\begin{verbatim}
python3 map_bvbrc.py              # 71/71 corpus map (BV-BRC)
python3 fetch_ncbi_proteomes.py   # 71 proteomes (NCBI Datasets)
python3 build_distribution.py     # BLASTP proteome comparison
python3 mge_context.py            # MGE/prophage + hotspot context
python3 crispr_survey.py          # CRISPR-Cas / RM known-system survey
python3 llm_judge.py              # LLM-judge verdict (Argo, free)
\end{verbatim}
Wall-clock $\sim$15 min (dominated by proteome downloads + 348k-protein BLAST DB).

\vspace{1em}
\noindent\textit{Report generated 2026-07-01, independent reproduction 2026-07-03. BVBRC-100 replication wave. Free endpoints only (BV-BRC, NCBI, Europe PMC, Argo).}

\end{document}
